\documentclass[11pt]{article}

\usepackage[a4paper,margin=1in]{geometry}
\usepackage{amsmath,amssymb,mathtools}
\usepackage{enumitem}
\usepackage{hyperref}

\newcommand{\uhat}{\widehat u}
\newcommand{\vh}{v_h}
\newcommand{\wh}{w_h}
\newcommand{\muhat}{\widehat\mu}
\newcommand{\qh}{q_h}
\newcommand{\Th}{\mathcal T_h}
\newcommand{\Eh}{\mathcal E_h}
\newcommand{\Ehi}{\mathcal E_h^\circ}
\newcommand{\R}{\mathcal R}
\newcommand{\Vh}{V_h}
\newcommand{\Qh}{\boldsymbol V_h}
\newcommand{\Mh}{M_h}
\newcommand{\dd}{\,\mathrm d}

\title{Residual Accounting in the Torsion-Initialized HDG Semilinear Newton Script}
\author{}
\date{}

\begin{document}
\maketitle

\section{Purpose}

This note records how the current Python script
\path{scripts/torsion_equilibrium/hdg/hdg_torsion_initialized_newton.py} computes and
uses residuals.  The goal is to make explicit why, after remeshing, the
reported Newton merit can stop decreasing even though the scalar/primal
equation is essentially solved.

The most important implementation fact is the following:
\begin{quote}
The Newton correction is driven by only the scalar/primal HDG residual
block, while the default merit \texttt{residual\_norm=euclid} is the
Euclidean norm of the full mixed HDG residual, including the flux
equations and the interior trace equation.
\end{quote}
Thus a defect introduced in the transferred flux or rebuilt trace can
remain visible in the merit even if the Newton step keeps reducing the
scalar equation for \(u_h\).

\section{Discrete Unknowns}

On each element \(K \in \Th\), the script stores
\[
  u_h|_K = \sum_i U_{K,i}\phi_i,
  \qquad
  q_{x,h}|_K = \sum_i Q^x_{K,i}\phi_i,
  \qquad
  q_{y,h}|_K = \sum_i Q^y_{K,i}\phi_i.
\]
The face trace unknown is stored per mesh edge:
\[
  \uhat_h|_e = \sum_a \widehat U_{e,a}\psi_a,
  \qquad e \in \Eh.
\]
Only interior edge residuals are included in the global trace residual;
boundary trace values are fixed by the homogeneous Dirichlet condition.

The nonlinear source is
\[
  \rho(u_h) =
  A\left(
    \sigma_{\varepsilon}(u_h-c_1)
    -
    \sigma_{\varepsilon}(u_h-c_2)
  \right),
\]
where \(A=\texttt{rho\_amp}\), and the derivative used in the Newton
linearization is
\[
  \rho'(u_h) =
  A\left(
    \sigma'_{\varepsilon}(u_h-c_1)
    -
    \sigma'_{\varepsilon}(u_h-c_2)
  \right).
\]
The quadrature values of \(\rho(u_h)\) are projected only when a
\texttt{DGField} representation is needed for metrics or plotting.  The
residual source moments are assembled directly from quadrature values.

\section{Weak HDG Residual Operators}

This section states the weak residuals before choosing a basis.  It is
the mathematical object that the coefficient residual in the script is
sampling.  The sign convention used by the HDG solver is
\[
  q \approx -\nabla u,
  \qquad
  -\nabla\cdot q = \rho(u).
\]
Equivalently, the scalar second order equation is
\(-\Delta u=\rho(u)\).

For each element \(K\), define the local scalar DG space
\[
  \Vh(K) = \mathbb P_p(K)
\]
and the local vector DG space
\[
  \Qh(K) = [\mathbb P_p(K)]^2.
\]
For each edge \(e\), define the trace space
\[
  \Mh(e)=\mathbb P_p(e).
\]
The global DG spaces are products over elements:
\[
  \Vh(\Th)=\prod_{K\in\Th}\Vh(K),
  \qquad
  \Qh(\Th)=\prod_{K\in\Th}\Qh(K),
\]
and the hybrid trace space is a product over edges.  The boundary trace
is prescribed by the Dirichlet data; in the present script it is zero.
Only the interior-edge trace equations are assembled into the residual.

Given an element state
\[
  (u_K,q_K,\widehat u_{\partial K})
  \in
  \Vh(K)\times\Qh(K)\times\prod_{e\subset\partial K}\Mh(e),
\]
the local HDG residual operator has two parts:
\[
  \R_K(u_K,q_K,\widehat u_{\partial K})
  =
  \left(
    \R^u_K(u_K,q_K,\widehat u_{\partial K}),
    \R^q_K(u_K,q_K,\widehat u_{\partial K})
  \right),
\]
with
\[
  \R^u_K(u_K,q_K,\widehat u_{\partial K})\in \Vh(K)',
  \qquad
  \R^q_K(u_K,q_K,\widehat u_{\partial K})\in \Qh(K)'.
\]

For every scalar test function \(v\in\Vh(K)\), the primal residual is
\[
  \boxed{
  \R^u_K(u_K,q_K,\widehat u_{\partial K};v)
  =
  -(q_K,\nabla v)_K
  +
  \left\langle
    q_K\cdot n_K+\tau(u_K-\widehat u_{\partial K}),
    v
  \right\rangle_{\partial K}
  -
  (\rho(u_K),v)_K .
  }
\]
This is the elementwise weak form of
\(-\nabla\cdot q=\rho(u)\), with the HDG numerical normal flux
\[
  \widehat f_K
  =
  q_K\cdot n_K+\tau(u_K-\widehat u_{\partial K}).
\]

For every vector test function \(r=(r_x,r_y)\in\Qh(K)\), the flux
residual is
\[
  \boxed{
  \R^q_K(u_K,q_K,\widehat u_{\partial K};r)
  =
  (u_K,\nabla\cdot r)_K
  -
  (q_K,r)_K
  -
  \left\langle
    \widehat u_{\partial K},r\cdot n_K
  \right\rangle_{\partial K}.
  }
\]
If the trace were exactly the element boundary value of \(u_K\), then
integration by parts would give
\[
  (u_K,\nabla\cdot r)_K
  -
  \langle u_K,r\cdot n_K\rangle_{\partial K}
  =
  -(\nabla u_K,r)_K,
\]
so \(\R^q_K=0\) enforces \(q_K=-\nabla u_K\).

The trace residual is not purely local to one element.  For each
interior edge \(e\in\Ehi\), define
\[
  \R^{\widehat u}_e(u_h,q_h,\widehat u_h)\in \Mh(e)'.
\]
For every trace test function \(\mu\in\Mh(e)\),
\[
  \boxed{
  \R^{\widehat u}_e(u_h,q_h,\widehat u_h;\mu)
  =
  \sum_{K:\,e\subset\partial K}
  \left\langle
    q_K\cdot n_K+\tau(u_K-\widehat u_h),
    \mu
  \right\rangle_e .
  }
\]
This is the weak conservation equation for the numerical normal flux on
the interior mesh skeleton.  There is no corresponding equation
\(\R^{\widehat u}_e=0\) on a physical boundary edge in this Dirichlet
formulation.  On boundary edges the trace \(\widehat u_h\) is prescribed
by the boundary condition, while the numerical boundary flux
\[
  q_K\cdot n_K+\tau(u_K-\widehat u_h)
\]
is generally nonzero and represents the boundary flux of the solution.
Consequently, boundary-edge flux is not included in \texttt{resTrace}.

The full residual operator is therefore
\[
  \R_{\mathrm{HDG}}(u_h,q_h,\widehat u_h)
  =
  \left(
    \{\R^u_K\}_{K\in\Th},
    \{\R^q_K\}_{K\in\Th},
    \{\R^{\widehat u}_e\}_{e\in\Ehi}
  \right).
\]
The vector returned by \texttt{hdg\_residual(...)} is obtained by
applying these functionals to the local basis functions:
\[
  (R^K_u)_i
  =
  \R^u_K(u_K,q_K,\widehat u_{\partial K};\phi_i),
\]
\[
  (R^K_{q_x})_i
  =
  \R^q_K(u_K,q_K,\widehat u_{\partial K};(\phi_i,0)),
  \qquad
  (R^K_{q_y})_i
  =
  \R^q_K(u_K,q_K,\widehat u_{\partial K};(0,\phi_i)),
\]
and
\[
  (R^e_{\widehat u})_a
  =
  \R^{\widehat u}_e(u_h,q_h,\widehat u_h;\psi_a).
\]
This is why the residual is best interpreted as a coefficient vector of
weak residual functionals.  With \texttt{residual\_norm=euclid}, the
script takes an ordinary Euclidean norm of this coefficient vector.  It
is not applying a mass-scaled dual norm unless \texttt{residual\_norm=hdg}
is selected.

\section{Element Matrices Used by the Residual}

Let \(\{\phi_i\}\) be the element basis and \(\{\psi_a\}\) the reference
edge basis.  The code uses affine element maps and precomputed local
matrices.  Up to the code's orientation conventions, the relevant
objects are:
\[
  M^K_{ij} = \int_K \phi_j\phi_i \dd x,
\]
\[
  D^{x,K}_{ij} = \int_K \phi_j\,\partial_x\phi_i \dd x,
  \qquad
  D^{y,K}_{ij} = \int_K \phi_j\,\partial_y\phi_i \dd x,
\]
This is the legacy-oriented derivative matrix used by the code: the
row index \(i\) belongs to the test basis function and the column index
\(j\) belongs to the trial coefficient.  With this orientation,
\(D^{x,K}Q^{x,K}\) represents \((q_{x,h},\partial_x\phi_i)_K\), not
\((\partial_x q_{x,h},\phi_i)_K\).
\[
  B^K_{\tau,ij}
  =
  \int_{\partial K} \tau\,\phi_j\phi_i \dd s,
\]
\[
  N^{x,K}_{ij}
  =
  \int_{\partial K} n_x\,\phi_j\phi_i \dd s,
  \qquad
  N^{y,K}_{ij}
  =
  \int_{\partial K} n_y\,\phi_j\phi_i \dd s.
\]
There is also a boundary coupling matrix \(E^K_\tau\) multiplying the
local trace coefficients on the three faces of \(K\).  In the code this
is \texttt{element\_boundary = diffusion\_element\_boundary\_mats(...)}.

The source moment vector on \(K\) is
\[
  F^K_i(u_h)
  =
  \int_K \rho(u_h)\phi_i \dd x.
\]
In code this is \texttt{scalar\_moments\_from\_values(space,
source\_values)}.

\section{Full Mixed HDG Residual}

The function \texttt{hdg\_residual(...)} returns one long vector
\[
  R_{\mathrm{mix}}
  =
  \begin{bmatrix}
    R^u \\
    R^q \\
    R^{\widehat u}
  \end{bmatrix}.
\]
The local element part has three blocks per element:
\[
  R^K_{\mathrm{loc}}
  =
  \begin{bmatrix}
    R^K_u \\
    R^K_{q_x} \\
    R^K_{q_y}
  \end{bmatrix}
  \in \mathbb R^{3N_p}.
\]

\subsection{Primal block}

The first block is the scalar/primal HDG equation:
\[
  R^K_u
  =
  B^K_\tau U^K
  + (N^{x,K}-D^{x,K})Q^{x,K}
  + (N^{y,K}-D^{y,K})Q^{y,K}
  - E^K_{\tau,u}\widehat U^K
  - F^K(U).
\]
This is the block reported as \texttt{resPrimal}.  It contains the
semilinear source \(\rho(u_h)\).

\subsection{Flux blocks}

The second and third local blocks enforce the first order relation
between \(q_h\) and \(u_h\):
\[
  R^K_{q_x}
  =
  D^{x,K}U^K
  - M^K Q^{x,K}
  - E^K_{\tau,q_x}\widehat U^K,
\]
\[
  R^K_{q_y}
  =
  D^{y,K}U^K
  - M^K Q^{y,K}
  - E^K_{\tau,q_y}\widehat U^K.
\]
These two blocks are grouped together and reported as \texttt{resFlux}.
They do not contain \(\rho(u_h)\).  They measure whether the stored flux
coefficients are consistent with the stored scalar field and trace on
the current mesh.

\subsection{Trace block}

For an interior edge \(e\), the script accumulates contributions from
the adjacent elements:
\[
  R^{e}_{\widehat u}
  =
  \sum_{K\supset e}
  \int_{e}
  \left(
    q_h^K\cdot n_K
    + \tau(u_h^K-\uhat_h)
  \right)\muhat
  \dd s .
\]
Only \(e\in\Ehi\) is retained in the returned residual vector.  This is
reported as \texttt{resTrace}.  It measures normal flux conservation and
the stabilization jump across the hybridized skeleton.

\section{Scalar Volume Residual}

The script also computes a diagnostic residual called
\texttt{resVolume}:
\[
  R^K_{\mathrm{vol},i}
  =
  \int_K \nabla u_h\cdot\nabla\phi_i\dd x
  -
  \int_K \rho(u_h)\phi_i\dd x.
\]
This is assembled by \texttt{scalar\_volume\_residual}.  It is useful for
comparison with the scalar continuous Galerkin or FreeFEM intuition, but
it is not the same as the full mixed HDG residual.  In particular,
\texttt{resVolume} ignores the stored \(q_h\) and \(\uhat_h\).

\section{Reported Norms}

After assembling \(R_{\mathrm{mix}}\), the script splits it as follows.
Let \(N_T\) be the number of elements and \(N_p\) the local scalar
polynomial dimension.  The local part has size \(3N_TN_p\).  The
remainder is the interior trace residual.

The reported coefficient norms are:
\[
  \texttt{resPrimal}
  =
  \|R^u\|_2,
\]
\[
  \texttt{resFlux}
  =
  \|R^q\|_2
  =
  \left\|
    \begin{bmatrix}
      R^{q_x}\\R^{q_y}
    \end{bmatrix}
  \right\|_2,
\]
\[
  \texttt{resTrace}
  =
  \|R^{\widehat u}\|_2,
\]
\[
  \texttt{resFull}
  =
  \|R_{\mathrm{mix}}\|_2.
\]
With the default command line used in the current runs,
\[
  \texttt{residual\_norm=euclid},
  \qquad
  \texttt{merit}=\texttt{resFull}.
\]

Other merit modes exist:
\begin{align*}
  \texttt{hdg-local}:&\quad \texttt{merit}=\texttt{resPrimal},\\
  \texttt{edp-volume}:&\quad \texttt{merit}=\texttt{resVolume},\\
  \texttt{hdg}:&\quad \texttt{merit}=\sqrt{R_{\mathrm{mix}}^T G^{-1}R_{\mathrm{mix}}},
\end{align*}
where \(G^{-1}\) is the HDG Gram inverse.  The Gram path is currently
not the default because of the separate Gram-matrix issue.

\section{Newton Correction Actually Solved}

At every Newton step the code first extracts only the primal block:
\[
  b^K_i = -R^K_{u,i}.
\]
This is exactly what \texttt{mixed\_u\_block\_rhs\_from\_residual}
returns.

Then the correction \((\delta u_h,\delta q_h,\delta\uhat_h)\) is
obtained by solving a linear HDG diffusion--reaction problem:
\[
  -\Delta(\delta u_h) - \rho'(u_h)\delta u_h
  =
  b,
\]
with the same HDG machinery.  In code:
\[
  \texttt{source\_h} \leftarrow b,
  \qquad
  \texttt{reaction\_h} \leftarrow -\rho'(u_h).
\]
The accepted trial state is then
\[
  u_h^{\mathrm{new}}
  =
  u_h+\alpha\delta u_h,
\]
\[
  q_h^{\mathrm{new}}
  =
  q_h+\alpha\delta q_h,
\]
\[
  \uhat_h^{\mathrm{new}}
  =
  \uhat_h+\alpha\delta\uhat_h.
\]

This explains the intended behavior on a fixed mesh.  If the current
triple \((u_h,q_h,\uhat_h)\) is already a reasonably consistent HDG
state, the correction generated by the scalar/primal block also carries
compatible flux and trace increments, so all blocks can decrease
together.

\section{What Changes After Remeshing}

During scheduled adaptation, the current implementation transfers the
state as follows:
\begin{enumerate}[label=(\arabic*)]
  \item transfer \(u_h\) to the new space;
  \item transfer \(q_{x,h}\) and \(q_{y,h}\) componentwise to the new
  space;
  \item rebuild \(\uhat_h\) by projecting the new element-face values of
  \(u_h\) to interior edge traces:
  \[
    \uhat_h
    =
    \operatorname{argmin}_{\widehat v_h}
    \sum_{K}\|u_h|_{\partial K}-\widehat v_h\|^2_{\partial K}.
  \]
\end{enumerate}
This is not the same as solving the mixed HDG equations on the new mesh.
It produces a geometrically reasonable scalar field and a geometrically
reasonable trace, but it does not guarantee
\[
  R^q = 0,
  \qquad
  R^{\widehat u}=0.
\]
The transferred flux is especially delicate because \(q_h\) should be
the HDG auxiliary variable associated with the current \(u_h\) and
\(\uhat_h\), not just the \(L^2\)-projection of the old flux field.

\section{Why the Current Plateau Occurs}

The recent diagnostic run showed the characteristic pattern
\[
  \texttt{resPrimal}\approx 10^{-9},
  \qquad
  \texttt{resFlux}\approx 6\times 10^{-5},
  \qquad
  \texttt{resTrace}\approx 1.5\times 10^{-4}.
\]
Since the default merit is
\[
  \texttt{merit}
  =
  \sqrt{
    \texttt{resPrimal}^2
    +
    \texttt{resFlux}^2
    +
    \texttt{resTrace}^2
  },
\]
the full merit cannot go below the flux/trace floor:
\[
  \texttt{merit}
  \approx
  \sqrt{
    (6\times 10^{-5})^2
    +
    (1.5\times 10^{-4})^2
  }
  \approx
  1.6\times 10^{-4}.
\]
This matches the observed plateau.

The Newton iteration keeps solving a correction whose right-hand side is
only \(-R^u\).  Once \(R^u\) is nearly zero, the next correction is very
small.  Therefore the update
\[
  (q_h,\uhat_h)
  \leftarrow
  (q_h,\uhat_h)+\alpha(\delta q_h,\delta\uhat_h)
\]
is also very small, and it has no strong mechanism left to remove a
pre-existing transfer defect in \(R^q\) or \(R^{\widehat u}\).  The line
search then sees no sufficient decrease in the full mixed merit and
eventually reports \texttt{FAIL\_LS}.

\section{Interpretation}

The current residual accounting says:
\begin{itemize}
  \item \texttt{resPrimal} is the semilinear scalar HDG equation being
  actively driven by the Newton source.
  \item \texttt{resFlux} is consistency of the stored auxiliary flux
  with \(u_h\) and \(\uhat_h\) on the current mesh.
  \item \texttt{resTrace} is interior skeleton conservation of
  \(q_h\cdot n+\tau(u_h-\uhat_h)\).
  \item \texttt{resFull} is the Euclidean coefficient norm of all of
  these blocks together.
\end{itemize}

On a fixed mesh, the HDG correction can reduce the full residual because
the state remains in the same discrete space.  After adaptation, the
state transfer can introduce flux/trace inconsistency.  Since the Newton
right-hand side is formed only from the primal block, the algorithm can
drive the scalar equation down while leaving the transferred auxiliary
defect almost unchanged.

\section{Implications for the Next Implementation Step}

The residual split suggests that the next experiment should not focus on
the scalar nonlinear source.  Instead, it should test how the auxiliary
state is rebuilt after adaptation.  Natural diagnostics are:
\begin{enumerate}[label=(\arabic*)]
  \item after transferring \(u_h\), reconstruct \(q_h\) and \(\uhat_h\)
  from an HDG local/global solve or projection consistent with the new
  mesh, instead of transferring old \(q_h\);
  \item compare the current transfer mode with a mode that discards
  transferred \(q_h\) and recomputes auxiliary variables from the new
  \(u_h\);
  \item optionally run the line search merit with \texttt{hdg-local}
  only as a diagnostic, not as a final correctness criterion, to confirm
  that the scalar equation is what is converging.
\end{enumerate}

\end{document}
